\documentclass[10pt,a4paper]{amsart}
\usepackage[margin=19mm]{geometry}
\usepackage{amsmath,amssymb,mathtools}
\usepackage[scr=rsfs,cal=euler]{mathalfa}
\usepackage{xfrac}
\usepackage[unicode,hidelinks,bookmarks=false]{hyperref}
\newcommand{\ie}{i.e.\ }
\newcommand{\cf}{cf.\ }
\newcommand{\resp}{respectively\ }
\newcommand{\Subsection}{\subsection}
\providecommand{\nobreakdash}{\nobreak\hbox{-}}
\setlength{\emergencystretch}{2em}
\multlinegap0pt
\usepackage{bm}
\usepackage{bbm}
\usepackage{dsfont}
\usepackage{xspace}
\usepackage{cancel}
\usepackage{mathtools}
\usepackage{amsthm}
\makeatletter
\@ifundefined{theorem}{\newtheorem{theorem}{Theorem}}{}
\@ifundefined{lemma}{\newtheorem{lemma}[theorem]{Lemma}}{}
\makeatother
\newcommand{\G}{\mathfrak G}
\newcommand{\Witt}{\mathrm{Witt}}
\newcommand{\ad}{\operatorname{ad}}
\newcommand{\epf}{\end{proof}}
\newcommand{\fB}{\mathfrak{B}}
\newcommand{\fG}{\mathfrak{G}}
\newcommand{\fM}{\mathfrak{M}}
\newcommand{\fb}{\mathfrak{b}}
\newcommand{\fg}{\mathfrak{g}}
\newcommand{\fm}{\mathfrak{m}}
\newcommand{\g}{\mathfrak g}
\newcommand{\pf}{\begin{proof}}
\renewcommand{\d}{\mathrm{d}}
\title[Weakly Noetherian Lie Algebra and the Sierra-Walton Conjectu]{Weakly Noetherian Lie Algebra and the Sierra-Walton Conjecture}
\author{Olivier Mathieu}
\date{Source: arXiv:2605.18116v1 (CC BY 4.0)}
\begin{document}
\maketitle

\noindent\emph{Source and licence.} Weakly Noetherian Lie Algebra and the Sierra-Walton Conjecture. Version arXiv:2605.18116v1 (CC BY 4.0), distributed under the Creative Commons Attribution 4.0 International licence, as recorded in the arXiv licence metadata for this exact version and shown on the abstract page https://arxiv.org/abs/2605.18116v1. Full rights notice: licenses/arxiv-2605.18116-CC-BY-4.0.txt.\par\smallskip

\noindent\textit{Editorial scope.} Counted unit: the Lemma whose source label is \texttt{Noeth-just} in the article (source0.tex lines 1409--1433, source label \texttt{Noeth-just}), delivered together with the article's own complete original proof of it (source0.tex lines 1435--1517, one \texttt{proof} environment of 83 lines). No mathematics is written, repaired or re-derived: statement and proof are the article's own text line for line. Required context retained uncounted: none. Every definition, display and label of those lines is retained unchanged. Omitted from this package and not redistributed: 1--1408, 1434--1434, 1518--4371. Nothing inside a retained excerpt is omitted or altered: every retained body is delivered exactly as the article's lines are written, so no omission receipt of any kind accompanies this package. Citations: the retained excerpts carry no citation command at all, so no bibliography entry of the article is transcribed and no reference is redistributed. Reference adaptation: none was needed and none was made; every reference inside the delivered unit resolves inside it, so no unresolved reference or citation is produced. Printed slips kept literal: no printed word, symbol, hypothesis or number of the counted statement or of its proof was altered. Layout and class: the article's own class is \texttt{amsart} with packages including amssymb, amsmath, amsthm, color, marginnote, hyperref, cancel, amscd, amsfonts, mathtools, enumitem, xy, fontenc, mathrsfs; the wrapper is the lane's pinned \texttt{amsart} editorial wrapper at 10pt on A4, which restates the theorem-like environments the retained text needs and adds the article's own macro definitions that the retained text needs (\texttt{\textbackslash G}, \texttt{\textbackslash Witt}, \texttt{\textbackslash ad}, \texttt{\textbackslash d}, \texttt{\textbackslash epf}, \texttt{\textbackslash fB}, \texttt{\textbackslash fG}, \texttt{\textbackslash fM}, \texttt{\textbackslash fb}, \texttt{\textbackslash fg}, \texttt{\textbackslash fm}, \texttt{\textbackslash g}, \texttt{\textbackslash pf}), with extra wrapper packages (bm, bbm, dsfont, xspace, cancel, mathtools). The delivered numbering is the unit's own. Assets: no retained span contains a figure environment, a graphics-inclusion command, a PDF-page inclusion, a table or a TikZ picture; no image file is copied into this package, the package declares no asset dependency and no image file is redistributed with it. Licence: version arXiv:2605.18116v1 of this article is distributed under the Creative Commons Attribution 4.0 International licence, as recorded in the arXiv licence metadata for this exact version and shown on the abstract page https://arxiv.org/abs/2605.18116v1. A full rights notice is delivered at licenses/arxiv-2605.18116-CC-BY-4.0.txt. Characters: every retained excerpt is pure ASCII, so the delivered engine, XeLaTeX with its default Latin Modern text font, reports no missing character.
\begin{lemma}\label{Noeth-just}
 Let $\g$ be a filtered Lie algebra, with associated graded algebra $\fG$. If

 \begin{itemize}
\item[(a1)]The Lie algebra $\fG_{\geq 1}$
is weakly Noetherian, and

\item[(a2)] the Lie algebra 
$\fg(0)$ is strongly Noetherian,
\end{itemize}

\noindent then the Lie algebra $\g$ is strongly  Noetherian.

\medbreak
Moreover, if 

 \begin{itemize}
\item[(b1)] All  graded ideals   of any graded subalgebra $\fB$ of $\G_{\geq 1}$ have 
finite codimension in $\fB$, and

\item[(b2)] any subalgebra of  $\fg(0)$ is just-infinite,
\end{itemize}

\noindent then all subalgebras of $\g$ are  just-infinite.
\end{lemma}

\begin{proof}
Any subalgebra $\fb$ of $\fg$ admits a filtration

$$0=\fb(-1)\subset\fb(0)\subset \fb(1)\subset\ldots$$ 

\noindent defined by $\fb(n)=\fb\cap \fg(n)$ for all
$n\geq 1$. Let $\fB$ be the associated graded 
Lie algebra.




We now prove the first assertion, namely
that $\fb$ is finitely generated. 
Since $\fB$ is a Lie subalgebra of $\fG$, Hypothesis (a1)
implies that

$$\dim \fB_{\geq 1}/[\fB_{\geq 1},\fB_{\geq 1}]
<\infty.$$


Since $\fB_{\geq 1}$ is positively graded, we deduce that $\fB_{\geq 1}$ is generated by a finite set
$y_1,\dots,y_m$ of homogenous elements of degree respectively
$d_1,\ldots,d_m$. Hence, there are  elements $x_1,\cdots,x_m\in\fb$ such that

\centerline{$x_i\in\fb(d_i)$ and $y_i=x_i\mod \fb(d_i-1)$
for all $i=1,\ldots,m$.}
Let $\fb'$ be the subalgebra generated by  $x_1,\cdots,x_m$.  It is clear that, as  a vector space we have

$$\fb=\fb'+\fb(0).$$ 

\noindent Both subalgebras  $\fb'$ and $\fb(0)$ are   finitely generated, therefore
$\fb$ is finitely generated, which proves the first assertion.

We now turn to the second assertion, that is that $\dim\fb/\fm<\infty$ for any nonzero ideal
$\fm$ of $\fb$. Let 

$$0=\fm(-1)\subset\fm(0)\subset \fm(1)\subset\ldots$$ 

\noindent be the filtration of $\fm$ previously defined and let 
$\fM$ be the associated graded Lie algebra.


It is obvious that $\fm(0)$ is an ideal of  $\fb(0)$ and it is clear that $\fM_{\geq 1}$ is an ideal
of $\fB_{\geq 1}$. Therefore
$\fb(0)/\fm(0)$ and $\fB_{\geq 1}/\fM_{\geq 1}$ are finite dimensional by hypotheses. We deduce that $\fb/\fm$  is finite dimensional, which proves the second assertion.
 \epf



\subsection{The Novikov-Krichever algebras are strongly Noetherian}


\begin{lemma}\label{graded-just}


Let $\fB$ be a graded subalgebra of
$\Witt_{\geq 1}(K)$.
\begin{itemize}
\item[(a)]
The Lie algebra  $\fB$  is finitely generated.
\item[(b)] Any nonzero graded ideal $\fm$ of $\fB$  has 
finite codimension in $\fB$.
\end{itemize}
\end{lemma}

\pf Recall that $\Witt_{\geq 1}(K)$ has basis 
$\{L_n:=t^{n+1}\d/\d t\mid n\geq 1\}$ and Lie bracket

\begin{align*}
[L_n,L_m]=(m-n) L_{n+m}
\end{align*}

\noindent
 We observe that $\ad(L_n)(L_m)\neq 0$, except if $n=m$.
 Hence, for any $n\geq 1$,  $\Witt_{\geq 1}$ is a $K[\ad(L_n)]$-module of finite type.

We can assume $\fB\neq 0$, thus  $\fB$  contains some basis element $L_n$. Hence
$\fB$  is a  $K[\ad(L_n)]$-module of finite type.  Therefore  $\fB$ is finitely generated, which completes the proof of Assertion (a).

Similarly,  $\fm$ 
contains  an homogenous element $L_m$. We have just proved that $\fB/[L_m,\fB]$ has finite dimension, therefore $\fB/\fm$ is finite dimensional.
\end{proof}

\end{document}
